\scanpage{7}
\mspage{35}
\deleted{
Alors puisque (5.1)(iii) est vrai pour $(g,\bar f,\bar G)$ et pour $(g',\pi,\bar G)$, les faisceaux $\deleted{\unsure{(}}R^q\bar f'_*\deleted{\unsure{)}}\inserted{(}\bar g^*\bar G\inserted{)}\simeq\unsure{g}^*(R^q\bar f_*\deleted{\unsure{)}}\inserted{(}\bar G\inserted{)})$ et $\deleted{\unsure{(}}R^q\pi'_*\deleted{\unsure{)}}\inserted{(}\bar g^*\bar G\inserted{)}\simeq g'^*(R^q\pi_*\deleted{\unsure{)}}\inserted{(}\bar G\inserted{)})$ sont nuls pour $q>0$. Il \deleted{s'ensuite}\inserted{résulte} de la suite spectrale de Leray
\[
E_2^{p,q}=\deleted{\unsure{(}}R^pf'_*\deleted{\unsure{)}}\inserted{(}(\deleted{\unsure{(}}R^q\pi'_*\deleted{\unsure{)}}\inserted{(}\bar g^*\bar G\inserted{)})\Longrightarrow\deleted{\unsure{(}}R^n\bar f'_*\deleted{\unsure{)}}\inserted{(}\bar G\inserted{)}
\]
qu'on a aussi $\deleted{\unsure{(}}R^pf'_*\deleted{\unsure{)}}\inserted{(}\pi'_*\bar g^*\bar G\inserted{)}=0$ pour $q>0$.

Or $\pi'_*\bar g^*\bar G\simeq g'^*(\pi_*\bar G)$, d'où $\deleted{\unsure{(}}R^pf'_*\deleted{\unsure{)}}\inserted{(}g'^*(\pi_*\bar G)\inserted{)}=0$ pour $q>0$. L'injection $\pi_*\bar F\longrightarrow\pi_*\bar G$ est donc un \deleted{\unsure{bon}} effacement des foncteurs $(R^pf'_*)g'^*$ dans la catégorie des faisceaux de torsion\inserted{(IX.1.2.(v))}, d'où le résultat.
}
\margin{$]\ ]$}

\uline{Lemme (8.4):} Il suffit de vérifier le théorème (5.1)(i) (resp. (i) et (ii), resp. (iii) ) pour $(g,f,F)$ avec $f$ \uline{projectif.}

Démonstration. Supposons le théorème (5.1)(i) (resp. $\ldots$) vrai chaque fois que $f$ est projectif, et soit $(g,f,F)$ arbitraire. D'après (6.5) avec $P$ la propriété vide, on peut supposer $S$ strictement local et $g:s_o\longrightarrow S$ l'inclusion du point fermé. Considérons la famille $\unsure{\mathfrak F|\mathcal F}$ des sous-préschémas fermés $h:Y\longrightarrow X$ tels que $Y$ contienne la fibre fermée $X_o$ de $X/S$ et tels que pour chaque $F$ et pour chaque $q$ \deleted{Applicable}\inserted{pertinent,} le morphisme
\[
H^q(X,F)\longrightarrow H^q(X,h_*h^*F)
\]
\mspage{6}
